\chapter{Calibration}
\label{ch:calibration}

Three calibrations run from the panel. None of them needs tools beyond a force
gauge.

\begin{center}\small
\begin{tabular}{@{}L{30mm}L{26mm}L{62mm}@{}}
\toprule
\textbf{Calibration} & \textbf{Started with} & \textbf{Corrects} \\
\midrule
Force        & \keycap{SETUP} & The difference between commanded and actual
                                bond force \\
Ultrasonic   & \keycap{TEST}  & Nothing --- it measures and reports \\
Tachometer   & Settings row 8 & Drift in the Z axis speed sensor \\
\bottomrule
\end{tabular}
\end{center}

\section{Force calibration}
\label{sec:force-cal}

The machine commands force in grams, but the coil's actual output drifts with
the lever, the tool and the mounting. This procedure measures the error once
and corrects every subsequent force command by it.

\subsection*{What you need}

A force gauge reading at least 100\,g, positioned so the tool bears on it the
way it bears on a workpiece.

\subsection*{Procedure}

\begin{enumerate}[leftmargin=*]
\item Set \menuitem{Setup Force} on the Settings page to the force you want to
  calibrate at. Choose a value near your normal bond force. The default is
  30\,g.
\item Position the gauge under the tool.
\item Press \keycap{SETUP}. The \keycap{SETUP} lamp lights and the machine
  waits.
\item \textbf{Press and hold the right mouse button.} The coil loads to
  \menuitem{Setup Force} and the head descends onto the gauge.
\item Read the gauge.
\item \textbf{Release the button.} The load is removed and the head retracts.
\item The Force Setup page opens automatically. Enter the reading with
  $\uparrow$ and $\downarrow$ or \keycap{$+$} and \keycap{$-$}, then press
  \keycap{ENTER}.
\item The machine confirms \msg{Force offset saved}.
\end{enumerate}

\begin{wbnote}
Hold the \textbf{right} button, not the left. The head descends for as long as
you hold it and stops when the gauge stops it --- the machine deliberately does
not wait to reach a target height, because your gauge is what ends the
descent.
\end{wbnote}

\begin{wbnote}
Press \keycap{ESC/DEL} instead of \keycap{ENTER} to leave the existing
correction untouched. Do that if the run did not go cleanly.
\end{wbnote}

\subsection*{What it does}

The machine stores the difference between what you measured and what it
commanded, and subtracts it from every force command thereafter. A zero force
command stays zero --- the correction can never turn ``coil off'' into a push.

\begin{wbnote}
Each run \emph{replaces} the stored correction, it does not add to it. Running
force calibration twice in a row is harmless.
\end{wbnote}

\subsection*{When to run it}

After changing the tool, after anything is disturbed on the bond lever, if bond
quality drifts without a recipe change, and as part of routine checks.

\section{Ultrasonic test}
\label{sec:us-test}

Press \keycap{TEST}. The machine sweeps for resonance, drives the transducer
briefly at \menuitem{Tail Power} until \menuitem{Tail Energy} has been
delivered, and reports:

\begin{lcdscreen}
\begin{lcdverb}
Freq:58.123 kHz
Q:12.3
Power:0.5 W
\end{lcdverb}
\end{lcdscreen}

The report stays for ten seconds. The axes do not move during the test.

\subsection*{What to look for}

\begin{description}[leftmargin=!,labelwidth=26mm,style=nextline,font=\normalfont\bfseries]
\item[Frequency] Must lie between \menuitem{Scan Start} and
  \menuitem{Scan Stop}, and comfortably inside them --- not at the edge. If it
  is near an edge, widen the window; if it is outside, the machine cannot find
  resonance at all and welds will fail with \msg{LOW BONDING POWER}.
\item[Q] Record it. The absolute number matters less than the trend. A Q that
  falls steadily points at a loosening horn, a loose tool screw, or a
  degrading transducer stack.
\item[Power] Should be close to \menuitem{Tail Power}. Much lower suggests the
  drive cannot reach the transducer.
\end{description}

\subsection*{Setting the scan window}

Run \keycap{TEST}, note the frequency, and set \menuitem{Scan Start} and
\menuitem{Scan Stop} roughly 2\,kHz either side of it. The defaults
(58--62\,kHz) suit a transducer resonating near 60\,kHz.

A wider window is more tolerant of drift but takes marginally longer to sweep
and gives a coarser fit for a given \menuitem{Scan Points}. If you widen the
window, consider raising \menuitem{Scan Points} to keep the resolution.

\section{Tachometer calibration}
\label{sec:tach-cal}

Corrects a small zero error in the Z axis speed sensor, which otherwise makes
the axis creep or hold imperfectly.

\subsection*{Procedure}

\begin{enumerate}[leftmargin=*]
\item Make sure nothing obstructs the bond head --- it will move.
\item Go to the Settings page (\keycap{SAVE} from the configuration list).
\item Move the \texttt{*} to \menuitem{Start Tach. Cal.} and press
  \keycap{ENTER}.
\item The head moves to 6\,mm and holds still while the machine samples for two
  seconds.
\item The correction is applied and saved immediately.
\end{enumerate}

\begin{wbnote}
Unlike force calibration, each run \emph{adds} to the stored correction rather
than replacing it. Running it two or three times converges on a better result;
that is the intended way to use it.
\end{wbnote}

\subsection*{When to run it}

If the head drifts when it should be holding still, if you see
\msg{Z POS. NOT SETTLED}, or as part of periodic maintenance.

\section{What is not operator-calibrated}

Sensor gains, the LVDT scaling and the underlying force-coil characterisation
are set at manufacture and are not adjustable from the panel. If you suspect
one of them, the machine needs service.
